\section*{Aula 1}

Você aprendeu em Mecânica Clássica a obter a trajetória de partículas ou
sistemas de partículas submetidas a forças conhecidas. Para uma partícula 
de massa $m$, você deve resolver 3 EDOs:
\begin{equation}
m \frac{d^2 {\bm r}}{dt^2} = {\bm F}({\bm r}, {\bm v}, t)
\end{equation}
onde ${\bm r}(t)$ é o vetor posição (em relação a uma origem 
previamente escolhida) e ${\bm F}({\bm r}, {\bm v}, t)$ é a força total que atua sobre a massa. Para que essa equação seja útil, você deve conhecer 
a dependência de ${\bm F}$ com a posição e a velocidade da partícula 
(a chamada {\it lei de força}). Algumas forças discutidas no curso de Mecânica Clássica (MC) possuíam leis de força bem definidas, como as 
{\bf forças de vínculo}, que aprendemos a eliminar no formalismo de Lagrange.

Do ponto de vista clássico e fundamental, apenas dois tipos de força atuam sobre
uma partícula, a \textbf{força gravitacional} e a \textbf{força eletromagnética}.
A primeira é devida a todos os corpos massivos nas vizinhanças da massa $m$. 
A segunda, que só se manifesta se a massa $m$ tiver uma carga $q$, é devida a
todas as cargas elétricas nas vizinhanças da massa $m$.

\begin{observacao}[]{}
    {\bf Observação:} O comentário sobre a necessidade da massa $m$ ter 
    carga só se aplica se estivermos falando de um partícula isolada. Um 
    sistema de partículas carregadas, mesmo que neutro (como um átomo), 
    experimenta forças eletromagnéticas.
\end{observacao}

Tanto a força gravitacional quanto a força eletromagnética são expressas em termos de campos:

\begin{subequations}
\begin{equation}
{\bm F}_{G}({\bm r}, t) = -m \nabla \phi({\bm r}, t)    
\end{equation}
\begin{equation}
{\bm F}_{EM}({\bm r}, {\bm v}, t) = q \left[ {\bm E}({\bm r}, t) + {\bm v} \times {\bm B}({\bm r}, t) \right]
\end{equation}
\end{subequations}
onde $\phi$ é o potencial gravitacional (veja o cap. 5 do Marion),
${\bm E}$ é o campo elétrico e ${\bm B}$ é o \textbf{campo magnético}.
Note que esses campos são avaliados na posição instantânea da partícula. A 
relação desses campos com as distribuições de massa e carga nas vizinhanças é:
\begin{equation}
    \nabla^2 \phi({\bm r}, t) = 4\pi G \rho({\bm r}, t)
\end{equation}
onde $G$ é a constante de gravitação e $\rho$ é a densidade de massa 
responsável pelo campo. As equações de Maxwell que relacionam os campos
eletromagnéticos com as distribuições de carga e corrente são:

\begin{subequations}
\begin{equation}
    \nabla \cdot {\bm E} = 4\pi \rho_e \quad \text{(Lei de Gauss)} 
\end{equation}
\begin{equation}
\nabla \times {\bm E} = -\frac{1}{c} \frac{\partial {\bm B}}{\partial t} \quad \text{(Lei de Faraday)} 
\end{equation}
\begin{equation}
\nabla \cdot {\bm B} = 0 \quad \text{(Ausência de monopolos magnéticos)} 
\end{equation}
\begin{equation}
\nabla \times {\bm B} = \frac{4\pi}{c} {\bm j} + \frac{1}{c} \frac{\partial 
{\bm E}}{\partial t} \quad \text{(Lei de Ampère-Maxwell)}
\end{equation}
\end{subequations}
Se você conhece a evolução temporal das distribuições de massa e carga, 
pode, em princípio, resolver essas equações para determinar $\phi$, 
${\bm E}$ e ${\bm B}$ e, dessa forma, escrever a EDO de Newton para a 
partícula de massa $m$.

Essas são as equações fundamentais da FÍSICA CLÁSSICA. Você aprenderá ${\bm F} = m{\bm a}$ no curso de Mecânica Clássica e as Eqs. de Maxwell no curso de Eletromagnetismo.

\begin{observacao}[]{}
{\br Observação:} A Mecânica Estatística, com seus conceitos de Temperatura,
Entropia, potenciais termodinâmicos etc., aplica-se a sistemas de 
MUITAS PARTÍCULAS interagindo por forças usuais (magnéticas, de contato etc.) 
e respondendo a essas forças da maneira como uma partícula isolada faria. A
diferença é que, ao invés de resolver a equação de Newton (no caso da Mecânica
Estatística clássica) ou de Schrödinger (no caso da Mecânica Estatística
Quântica) para cada partícula, resolvem-se equações para as {DISTRIBUIÇÕES DE 
PROBABILIDADE} do sistema (isto é, introduzem-se conceitos estatísticos para
descrever as propriedades do sistema). De fato, no curso de Mecânica 
Estatística, nos concentramos mais nas propriedades das distribuições 
\textbf{INDEPENDENTES DO TEMPO}, isto é, no \textbf{EQUILÍBRIO TERMODINÂMICO}.
\end{observacao}
Essas equações da Física Clássica descrevem uma variedade tão grande de 
fenômenos e de maneira tão compacta e elegante que causou muito desconforto a 
observação de pequenas falhas nessas equações.

\begin{enumerate}
\item As equações de Maxwell prevêem que a velocidade da luz deve ser a mesma para todos os observadores inerciais, e isso é incompatível com a maneira como o tempo é tratado em ${\bm F} = m{\bm a}$ (\textbf{Teoria da Relatividade Especial} - 1905 - Lorentz, Poincaré, Einstein).

text
\item A equação da Gravitação de Newton diz que uma massa em qualquer ponto do universo influencia \textbf{instantaneamente} o movimento de qualquer outra massa em qualquer outro ponto do universo, e isso é incompatível com a Teoria da Relatividade Especial $\rightarrow$ \textbf{Teoria da Relatividade Geral} (1915 - Einstein).

\item O grande desastre da Física Clássica ocorreu quando se tentou explicar o comportamento de objetos muito pequenos, os átomos e as moléculas, e sua interação com a radiação eletromagnética.
\end{enumerate}

\subsection*{Histórico do Desenvolvimento da Teoria Quântica:}
\begin{itemize}
\item \textbf{1834} - Michael Faraday - Eletrólise de várias substâncias.
\item \textbf{1867} - Fórmulas para moléculas $\text{H}_2\text{O}$, $\text{H}_2$, $\text{CO}$.
\item \textbf{1897} - J.J. Thomson - Descoberta do elétron.
\item \textbf{1909} - R.A. Millikan - Medida da carga do elétron ($e$) e consequências:
\begin{itemize}
\item Valor de $m_e$ é obtido (dos resultados de Thomson)
\item Valor de $N_A = 6.022 \times 10^{23}$ é obtido (dos resultados de Faraday)
\item Tamanho do átomo $\sim 1$ Å é obtido (das densidades)
\item Massa dos átomos é obtida (de $N_A$)
\end{itemize}
\end{itemize}

A imagem que surge é a do átomo como uma pequena esfera composta de uma carga 
positiva massiva e de uma carga negativa muito leve (a massa dos elétrons é 
muito menor que a massa dos átomos). Como se dá o movimento e a distribuição 
dessas cargas no átomo? Esse movimento deve ser a causa do espectro 
característico das substâncias no estado gasoso (o movimento de cargas causa a 
emissão de radiação eletromagnética com frequência igual à frequência do 
movimento das cargas, segundo Maxwell).

\subsection*{Espectros Atômicos:}
J.R. Rydberg (1900) e W. Ritz (1908) mostraram que os comprimentos de onda 
característicos de cada substância podiam ser compreendidos como diferenças 
entre \textbf{termos}. Por exemplo, na série de Balmer do Hidrogênio:

\begin{equation}
\frac{1}{\lambda} = R_H \left( \frac{1}{2^2} - \frac{1}{n^2} \right) \quad (n = 
3, 4, 5, \dots)
\end{equation}
onde $R_H$ é a \textbf{constante de Rydberg} para o hidrogênio. O número 
de termos necessários para explicar um dado espectro era muito menor
que o número de comprimentos de onda característicos.

\subsection*{Modelo de Rutherford (1911):}
Nesse modelo, o núcleo, de tamanho $\sim 10^{-4}$ Å (muito menor que o átomo, 
de tamanho $\approx 1$ Å), faz o papel do Sol, e os elétrons, o dos planetas. 
O número de elétrons (e a carga do núcleo), o número atômico $Z$, foi medido 
por Rutherford.

O problema é que cargas orbitando um núcleo carregado deveriam emitir energia
eletromagnética continuamente até colapsar no interior do núcleo. Esse é o 
problema da \textbf{ESTABILIDADE DA MATÉRIA}, que ${\bm F} = m{\bm a}$ + Eqs. 
de Maxwell \textbf{não conseguem resolver}!

\section*{Resumo dos Problemas com ${\bm F} = m{\bm a}$ e Eqs. de Maxwell}
\begin{itemize}
\item \textbf{Átomo}:
\begin{itemize}
\item Estabilidade da matéria
\item Espectro discreto dos átomos dado como diferença de energias
\end{itemize}

text
\item \textbf{Ondas Eletromagnéticas}:
\begin{itemize}
    \item Energia de um campo EM clássico confinado em uma cavidade seria infinita (\textbf{Problema da Radiação de Corpo Negro} - Max Planck, 1900)
    \item Efeito fotoelétrico (Einstein, 1905)
    \item Espalhamento Compton (A.H. Compton, 1919-1923)
\end{itemize}
\end{itemize}

Nesses fenômenos, a radiação eletromagnética se comporta como partícula de energia $E = h\nu$ e momento $p = \frac{h}{\lambda}$ (onde $\nu$ e $\lambda$ são a frequência e o comprimento de onda da onda eletromagnética).

\subsection*{Espalhamento de Elétrons:}
Elétrons espalhados por uma superfície metálica apresentam um padrão de difração (a rede de difração são os átomos na superfície) - C.J. Davisson e L.H. Germer (1925).

Nesse experimento, o comprimento de onda da "onda eletrônica" é dado por:

\begin{equation}
\lambda = \frac{h}{p}
\end{equation}

onde $p$ é o momento usual do elétron clássico.

G.P. Thomson (1927-28) fez experiências semelhantes e provou, com um campo magnético, que o padrão de difração era devido a elétrons.

\subsection*{Conclusão:}
Vemos portanto que não apenas as equações da Física Clássica não conseguem explicar o comportamento interno dos átomos, mas os próprios conceitos de \textbf{PARTÍCULA} e \textbf{ONDA} parecem inadequados para descrever o universo na escala de distâncias atômicas ($L \sim 1$ Å).

\subsection*{O que acontece de especial nessa escala de tamanho?}
A Ciência depende de observações: medimos ${\bm r}(t)$ e verificamos que ${\bm F} = m{\bm a}$; medimos ${\bm E}$ e ${\bm B}$ e verificamos as Eqs. de Maxwell. A Física clássica pressupõe que o ato da observação pode ser feito de forma arbitrariamente não perturbadora. No caso do átomo, imaginamos que seria possível medir ${\bm r}(t)$ do elétron com uma radiação EM "muito delicada" (que não carregaria momento linear apreciável), de modo a não alterar o movimento eletrônico. Se isso fosse possível, poderíamos verificar ${\bm F} = m{\bm a}$, onde ${\bm F}$ seria a força elétrica do núcleo.

\textbf{MAS ISSO NÃO É POSSÍVEL!} Não há como medir a posição do elétron sem alterar sua velocidade, e isso não é uma impossibilidade técnica, mas sim uma impossibilidade de fato.

\begin{quote}
"A NATUREZA TEM UM LIMITE FUNDAMENTAL PARA O PEQUENO: aquele que, ao ser observado, é sempre perturbado, independente do esforço do observador."
\end{quote}

A descrição das partículas subatômicas não pode ser em termos de trajetórias. É aqui que entra a \textbf{Mecânica Quântica} e sua \textbf{onda de probabilidade}.

A razão da descrição ter necessariamente que ser probabilística está no fato de que, após cada medida, o sistema se modifica de maneira a tornar impossível prever o resultado de uma medida subsequente (exceto em alguns casos especiais). Tudo o que poderemos fornecer, em geral, é a \textbf{probabilidade} dos resultados das medidas.

Neste curso, trataremos da descrição quântica de partículas materiais. A descrição quântica do campo EM e o fóton terão que esperar.